% ====================================================================
% primer_body.tex -- plain-language story for a non-expert reader.
% Inputs only into main_primer.tex.  Numbers come from values.tex; figures
% are shared with the collaboration note via the common \graphicspath.
% ====================================================================

\section{What was measured, and why it is hard}

A neutrino is a nearly invisible particle: it passes through ordinary matter
almost without a trace. To study neutrinos, the MINERvA experiment at Fermilab
fired an intense beam of them at a finely segmented detector and recorded the
rare occasions when one struck a nucleus and produced visible particles --- in
this analysis, an interaction that creates a muon (a heavy cousin of the
electron). The quantity of interest is the \emph{cross section}: essentially the
probability that such an interaction happens, measured as a function of the
muon's motion --- how fast it travels along the beam ($\pPar$) and how much it is
deflected sideways ($\pT$).

The difficulty is that the detector never measures these quantities perfectly.
A muon's recorded momentum is a blurred version of its true momentum, the blur
being the detector's finite resolution and efficiency. The job of
\emph{unfolding} is to remove that blur: to recover the true distribution that
nature produced from the smeared distribution the detector recorded. It is the
same problem as deblurring a photograph, and like deblurring it is delicate ---
done carelessly it amplifies noise into artifacts.

\section{The method, in plain language}

The traditional way to unfold builds a big table --- a ``migration matrix'' ---
that records, for every bin of true momentum, how events leak into neighboring
bins of measured momentum, and then tries to invert that leakage. It works, and
nothing forbids it in principle; the limitation is one of scale. Each new
quantity you want to measure multiplies the number of bins, so the simulation
needed to fill the table --- and the table's own conditioning --- degrade fast.
In published practice that has kept binned results at two, occasionally three,
reported quantities at a time.

This analysis uses a different, machine-learning method called \textbf{OmniFold}.
Instead of a table of bins, it trains a computer to compare the real data with a
detailed computer simulation, event by event, and to learn a \emph{weight} for
each simulated event --- a number that says ``there should be a little more (or
less) of this kind of event to make the simulation look like the data.'' Because
the simulation knows both the true and the measured version of every event, that
weight can be carried back from the measured world to the true world. Repeating
the comparison a few times morphs the simulation until it matches the data, and
the morphed \emph{true} simulation is the unfolded answer.

The decisive advantage is what happens when you want to measure more quantities
at once. In the table method, adding a quantity multiplies the bins. In
OmniFold, adding a quantity adds an input feature without requiring a new table
of migration bins. Training, the amount of simulation needed, and uncertainty
estimation still become more demanding. This makes the extensions in the next
section practical without removing the need to validate them.

\begin{figure}[H]
  \centering
  \includegraphics[width=0.8\textwidth]{migration_resolution}
  \caption{Why unfolding is needed: the detector mixes neighboring values of
  the true muon momentum together (the off-diagonal spread in these maps). A
  na\"ive bin-by-bin correction would be badly biased; unfolding undoes the
  mixing. OmniFold does it without ever building these bins.}
\end{figure}

\section{What we found}

The multi-quantity central values below pass the stated validation tests. A
five-quantity uncertainty has now been adopted, and the two-quantity
$(\Eavail,W)$ version of it has been built --- but the multi-quantity plots
still show central values only, and we still do not assign statistical
significance to their differences from the simulations. The reason is worth
stating plainly, because it is something we measured rather than a precaution.

Our method has a random component: the learner is started from a random seed.
Change that seed, leaving everything else alone, and the \emph{uncertainty} we
report on the projected result shifts by about \SI{\sprojMeasured}{\percent}
--- more than the \SI{\sprojBound}{\percent} we had decided in advance we
would tolerate. The natural hope is that this is just statistical noise from
having a finite number of simulated variations, which would shrink if we ran
more of them. We checked over the range we have, and it does not shrink: the
shift is the same size whether we use a quarter, a half, or all of our
variations, while genuine statistical noise, measured between the quarter and
the half, falls by a factor of nearly three over the same step. So the shift
behaves like a real property of the learner rather than like noise, and we saw
no reduction between a quarter and the full set. What we did \emph{not}
establish is what a much larger set would give: we compared only one pair of
random seeds, and nested subsets of one set of variations cannot rule out every
larger one. Two further facts qualify it: five of the forty-five uncertainty
components are
built in a way the seed cannot reach at all, so the figure above is the effect
of the other forty and we have not measured which way those five would move it;
and the same seed change moves the \emph{central values} too,
though only by at most \SI{\cvMoveOverSigmaMax}{\percent} of their own
uncertainty. Reporting a significance on top of that would claim a precision
the input does not have.

One more thing is worth saying plainly about what the detector uncertainties
cover. Five of them describe effects that move a muon's measured momentum or
angle, and for those we rebuilt the simulation end to end so that events
crossing into or out of the selection are counted properly. Four others ---
the MINOS efficiency and three describing how hadrons interact in the detector
--- are applied instead by reweighting the simulated events, so there is
nothing to rebuild. That tells us why those four sit outside the rebuild; it
does not by itself show that the set of hadronic-response variations we
include is the complete set this measurement needs.

The two-dimensional trust check is unaffected and retains its independently
constructed uncertainty.

\subsection*{First, a trust check: reproduce a known result}

Before trusting a new method, one must show it reproduces an established
measurement. MINERvA already published this muon cross section using the
traditional method. Our machine-learning result agrees with it closely: the
overall interaction rate matches to about one percent (ratio $\ratioTot$), and
\SI{\binsTen}{\percent} of the individual measurement cells agree within ten
percent. Just as important, our \emph{uncertainties} --- worked out completely
independently --- come out essentially the same size as the published ones
(\SI{\uqMedian}{\percent} versus \SI{\uqPaper}{\percent} typical). The new method
reproduces the old answer and its error bars.

\begin{figure}[H]
  \centering
  \includegraphics[width=0.62\textwidth]{MEFHC_5iter_fig13}
  \caption{The unfolded muon cross section as a function of the two momentum
  components. This reproduces a published MINERvA measurement, but was produced
  by the machine-learning method described above.}
\end{figure}

\subsection*{Then, measure more at once --- and rediscover known physics}

With the trust check passed, we used OmniFold's defining ability and unfolded
several quantities \emph{simultaneously} and \emph{without binning} --- a first
at MINERvA, where earlier multi-quantity results used binned methods: the muon
momentum together with the available energy ($\Eavail$, roughly how much energy
the struck nucleus released), the momentum transfer ($q_{3}$), and the hadronic
mass ($W$). Each richer measurement, when collapsed back down, reproduces the
simpler one --- a built-in consistency check.

The first physics payoff appears in the available energy. When we compare the
data to four state-of-the-art simulation predictions, all four fall short of
the data in the low-energy region (Fig.~\ref{fig:primergen}, left). This is not a
surprise --- it resembles an established low-recoil discrepancy that simulations
have historically struggled to model. Adding a process in which the neutrino
interacts with a pair of nucleons (``2p2h'') fills part of the gap. The two tested
pion-rescattering variations have a much smaller effect, but that does not rule
out other changes to how particles move through the nucleus. Related MINERvA
measurements with a different event selection show why proton and pion
rescattering also matter~\cite{MINERvA:2026qeleme}. These comparisons support the
two-nucleon contribution without uniquely establishing the cause. And
where MINERvA has a dedicated published measurement of this region, ours agrees
with it at central-value level on the maximal common grid.

\begin{figure}[H]
  \centering
  \includegraphics[width=0.99\textwidth]{eavailW_band}
  \caption{The unfolded central values (black points) compared to four tested
  simulation predictions --- GENIE's own central configuration, the same GENIE
  with the two-nucleon process added, NuWro and GiBUU. \emph{Left:} available
  energy; all four fall below the data at low available energy, the known
  ``low-recoil'' shortfall, here rediscovered by a different unfolding method.
  \emph{Right:} hadronic mass; all four sit below the data at high hadronic
  mass. \textbf{Central values only.} No uncertainty band is drawn and no
  significance is assigned, for the reason given at the start of this section.
  The first two predictions are two configurations of the same generator and
  are not independent of each other.}
  \label{fig:primergen}
\end{figure}

\subsection*{A new puzzle that no simulation explains}

Adding the hadronic-mass axis $W$ lets us look where the simulations still miss
\emph{after} the two-nucleon process is included. The leftover discrepancy sits
in a specific corner: high available energy \emph{and} high $W$ --- the
``deep-inelastic'' region where the neutrino shatters the nucleon. That corner
is a statement about the two quantities \emph{jointly}, so
Fig.~\ref{fig:primerjoint} shows the joint map rather than the two separate
one-dimensional views of Fig.~\ref{fig:primergen}: neither of those alone
would show it. There, all four simulations undershoot the central result. Its
statistical significance is not assigned. The random seed described at the start of this section moves
these central values as well, by at most \SI{\cvMoveOverSigmaMax}{\percent} of
their own uncertainty; how far the corner's share of the excess could shift
with it was not separately measured. This is the kind of localized discrepancy that
points theorists toward what their models are missing --- and it was only
visible because the unbinned method let us slice the data this finely.

\begin{figure}[H]
  \centering
  \includegraphics[width=0.95\textwidth]{paper_joint_localization}
  \caption{The unfolded data compared to the GENIE central simulation across
  the available-energy / hadronic-mass plane \emph{jointly}. Both panels are
  measured on the same seven-by-six grid of cells, with available-energy
  boundaries $0$, $0.1$, $0.2$, $0.4$, $0.8$, $1.5$, $3$ and $100$~GeV and
  hadronic-mass boundaries $0$, $1.1$, $1.4$, $1.8$, $2.2$, $3$ and
  $100$~GeV; the axes are labelled by cell number, counting from zero, and the
  last cell on each axis is a wide catch-all that collects everything above
  $3$~GeV. \emph{Left:} the ratio of data to simulation, so $1$ means
  agreement. \emph{Right:} the difference, data minus simulation, as a
  cross section \emph{integrated over each cell}
  (\si{cm^2/nucleon}) --- because that includes the cell's area, the wide
  catch-all cells look large here even where the ratio on the left is
  moderate, and the two panels should be read together. The excess concentrates
  in the top-right corner: high available energy \emph{and} high hadronic mass.
  Central values only; no uncertainty is drawn and no significance is assigned.
  The estimator seed moves these central values by at most
  \SI{\cvMoveOverSigmaMax}{\percent} of their own uncertainty.}
  \label{fig:primerjoint}
\end{figure}

\subsection*{Alternative representations and signal definitions}

Two studies test the method's limits rather than adding to the measurement.

We ran a ``point-cloud transformer'' on the raw detector deposits instead of
the summary quantities. It shows that the machinery runs end to end on low-level
detector data, pairing reconstructed recoil clusters with truth-side final-state
hadrons. The muon was withheld from both classifiers, so this is a recoil-only
method-development cross-check, not a measurement of the whole event; it
contributes no cross section or error bar. A full-event version would supply the
muon and test its added information directly.

Separately, the published signal definition excludes about a third of the
simulated interaction rate because of its muon kinematic cuts. This fraction
is a simulation estimate, not a measured number. We relaxed those truth-level cuts while retaining the
same tracker interaction volume and reconstructed MINOS-based selection. The
preliminary central value is
$\sigma_{\mathrm{ext}}=\SI{\fpsSigma}{cm^2/nucleon}$, but only about
\SI{6}{\percent} of the rate is a newly detector-supported extension; about
\SI{28}{\percent} lies in low- or zero-efficiency cells and remains a
separately labeled, simulation-dependent extrapolation.

\begin{figure}[H]
  \centering
  \includegraphics[width=0.99\textwidth]{fps_pilot_compare_MEFHC}
  \caption{The extended-fiducial scalar study, on the muon momentum grid
  ($p_T$ versus $p_\parallel$, in bin number). Left: inside the published region
  (between the black lines) the new result matches the original to better than
  one percent. Center: changing the starting simulation has a small effect in
  accepted cells but a much larger one in low-efficiency cells. Right: the
  preliminary central value over the enlarged truth domain (color = rate, on a
  logarithmic scale); detector-supported and prior-dominated regions must be
  interpreted separately.}
\end{figure}

\section{The takeaway}

Unbinned machine-learning unfolding reproduces a flagship MINERvA measurement
\emph{and} scales to many quantities at once, where the binned method's
response table does not. Using that, this analysis recovered a known low-recoil discrepancy with
a different unfolding method, exposed a localized central-value difference that all four tested
simulations underpredict, and
mapped a possible extension of the signal definition while exposing where the
detector provides too little information for a model-independent result.
